% !PRINTED 471
\weberpage{471}
Setzen wir nun in (13)

\[
\varDelta = (g\,C^3)^{\frac{1}{7}}, \quad f = h\,C^2\,(g\,C^3)^{-\frac{4}{7}},
\]

so ergiebt sich für den Ausdruck (13) mit Benutzung von (14)

\[
g^{b+2c}\,h^a\,C^{\frac{m}{14}}\,g^{-\frac{m}{7}},
\]

und im Zähler und Nenner lässt sich der Factor $C^{\frac{m}{14}}
g^{-\frac{m}{7}}$ heben, so dass alles rational durch $g$ und $h$
ausgedrückt ist.

Ebenso wie $g, h$ hängt auch die Function $u$ nur von dem Verhältniss
der Variablen $x_1 : x_2 : x_3$ ab, und es folgt leicht, dass jede
andere Function von derselben Eigenschaft, die wie $u$ die
Substitutionen der Gruppe $G_{24}$ gestattet, rational durch $u, g,
h$ darstellbar ist. Denn eine solche Function lässt sich zunächst
nach den allgemeinen Sätzen des \S.\,43 als rationale Function von
$u$ und den Invarianten darstellen, und zwar nur auf eine Weise als
ganze Function von $u$, die den $6^{\text{ten}}$ Grad nicht
übersteigt (Bd.\,I, \S.\,142). Die Co\"efficienten in dieser
Darstellung sind Invarianten, und da $u$ nur von den Verhältnissen
abhängt, so können auch die Co\"efficienten nur von den
Verhältnissen abhängen, und sind daher rational durch $g, h$
darstellbar.

\begin{center}
\S.\ 126.

\gesperrt{Reduction der allgemeinen Resolvente siebenten Grades auf
die specielle.}
\end{center}

Wir haben im vorigen Paragraphen zwei Formen der Resolvente
$7^{\text{ten}}$ Grades des Formenproblems kennen gelernt, von denen
die eine, die specielle, für den Fall gilt, dass $f = 0$ ist, d.\,h.
für den Fall, dass der gesuchte Punkt auf der Grundcurve liegt, die
allgemeine für den Fall, dass er eine beliebige Lage hat.

Die Grössen $u, g, h$ wollen wir, wenn sie sich auf den Punkt $(x)$
beziehen, mit

\weq{1}{u_x = \dfrac{\varDelta^2 z}{C}, \quad g_x = \dfrac{\varDelta^7}{C^3}, \quad h_x = \dfrac{f\varDelta^4}{C^2}}

bezeichnen. Die allgemeine Resolvente soll dann mit

\weq{2}{R(u_x, g_x, h_x) = R_x = 0}

bezeichnet werden, und die specielle geht daraus hervor, wenn
%CONT>>
